\begin{question}
\noindent $ABC$ is a straight line. $D$ is a point above the line, with rays drawn from $D$ to points $A$, $E$, $C$ where $E$ lies on $AC$ between $A$ and $C$.

The three angles formed at point on the line by the rays from $D$ are algebraic expressions, as shown in the diagram.

\begin{center}
\begin{tikzpicture}[scale=1.4]
    \coordinate (A) at (0,0);
    \coordinate (C) at (6,0);
    \coordinate (P) at (2.4,0);
    \coordinate (D) at (2.9,2.6);

    \draw[thick] (A) -- (C);
    \draw[thick] (A) -- (D);
    \draw[thick] (P) -- (D);
    \draw[thick] (C) -- (D);

    \node at (A) [below] {$A$};
    \node at (C) [below] {$C$};
    \node at (P) [below] {$B$};
    \node at (D) [above] {$D$};

    \pic [draw, "\footnotesize $2x^{\circ}$", angle eccentricity=1.7, angle radius=0.9cm] {angle = A--D--P};
    \pic [draw, "\footnotesize $(3x+10)^{\circ}$", angle eccentricity=1.6, angle radius=1cm] {angle = P--D--C};

    \pic [draw, "\footnotesize $(x+20)^{\circ}$", angle eccentricity=2.6, angle radius=0.5cm] {angle = D--A--P};
    \pic [draw, "\footnotesize $x^{\circ}$", angle eccentricity=2.2, angle radius=0.5cm] {angle = D--C--P};

\end{tikzpicture}
\end{center}

\noindent In triangle $ABD$, the angle at $A$ is $(x+20)^{\circ}$ and angle $ADB = 2x^{\circ}$.\\
In triangle $CBD$, the angle at $C$ is $x^{\circ}$ and angle $BDC = (3x+10)^{\circ}$.\\
Angle $ABD$ and angle $DBC$ lie on the straight line $ABC$.

Since angles $ABD$ and $DBC$ are angles on a straight line, and using the fact that angles in each triangle sum to $180^{\circ}$, it can be shown that:
\[
2x + (3x+10) = 180
\]

Use this equation to work out the value of $x$.

\vspace{4cm}

\begin{flushright}
\answereq{x}{3}
\end{flushright}
\end{question}
